\documentclass[12pt]{article}
\usepackage{amsmath,amssymb,amsfonts}
\usepackage[margin=1in]{geometry}
\usepackage{bm}

\newcommand{\vect}[1]{\mathbf{#1}}
\newcommand{\ehat}[1]{\hat{\vect{e}}_{#1}}

\begin{document}

\section*{Problem 4}

\textbf{Problem:} A charged particle moves in a region of space where there is a uniform magnetic field $\vect{B}$ in the $x$-direction and a uniform electric field $\vect{E}$ in the $yz$-plane. Write the equations of motion for the particle. Then solve them, taking for initial conditions (at $t=0$): $x = y = z = 0$; $\dot{x} = x_0$, $\dot{y} = \dot{z} = 0$. Show that there is a drift velocity equal to $E_z/B$ in the positive $y$-direction. If $E_z = 0$, then the orbit becomes a spiral about a line parallel to the $z$-axis.

\bigskip
\textbf{Solution:}

\subsection*{Setting up the equations of motion}

We have $\vect{B} = B\,\ehat{x}$ and $\vect{E} = E_y\,\ehat{y} + E_z\,\ehat{z}$.

The Lorentz force equation (in Gaussian units) is:
\[
m\ddot{\vect{r}} = q\vect{E} + \frac{q}{c}\,\dot{\vect{r}} \times \vect{B}.
\]

Compute $\dot{\vect{r}} \times \vect{B}$:
\[
\dot{\vect{r}} \times \vect{B} = (\dot{x}\,\ehat{x} + \dot{y}\,\ehat{y} + \dot{z}\,\ehat{z}) \times B\,\ehat{x}
= B\dot{y}\,(\ehat{y}\times\ehat{x}) + B\dot{z}\,(\ehat{z}\times\ehat{x})
= -B\dot{y}\,\ehat{z} + B\dot{z}\,\ehat{y}.
\]

Defining the cyclotron frequency $\omega_c = qB/(mc)$, the component equations are:
\begin{align}
m\ddot{x} &= 0, \label{eq:xeom} \\
m\ddot{y} &= qE_y + \frac{qB}{c}\dot{z} = qE_y + m\omega_c \dot{z}, \label{eq:yeom} \\
m\ddot{z} &= qE_z - \frac{qB}{c}\dot{y} = qE_z - m\omega_c \dot{y}. \label{eq:zeom}
\end{align}

\subsection*{Solving the $x$-equation}

From Eq.~\eqref{eq:xeom}: $\ddot{x} = 0$, so with $x(0) = 0$, $\dot{x}(0) = x_0$:
\[
\boxed{x(t) = x_0 t}.
\]
The motion in the $x$-direction is uniform.

\subsection*{Solving the coupled $y$-$z$ equations}

Let $\alpha_y = qE_y/m$ and $\alpha_z = qE_z/m$. From Eqs.~\eqref{eq:yeom}--\eqref{eq:zeom}:
\begin{align}
\ddot{y} &= \alpha_y + \omega_c \dot{z}, \label{eq:y2} \\
\ddot{z} &= \alpha_z - \omega_c \dot{y}. \label{eq:z2}
\end{align}

Define $w = \dot{y} + i\dot{z}$. Then:
\[
\dot{w} = \ddot{y} + i\ddot{z} = (\alpha_y + \omega_c \dot{z}) + i(\alpha_z - \omega_c \dot{y})
= (\alpha_y + i\alpha_z) - i\omega_c(\dot{y} + i\dot{z})
= (\alpha_y + i\alpha_z) - i\omega_c w.
\]

This is a first-order ODE: $\dot{w} + i\omega_c w = \alpha_y + i\alpha_z$.

The solution with $w(0) = 0$ (since $\dot{y}(0) = \dot{z}(0) = 0$) is:
\[
w(t) = \frac{\alpha_y + i\alpha_z}{i\omega_c}\bigl(1 - e^{-i\omega_c t}\bigr).
\]

Computing the constant:
\[
\frac{\alpha_y + i\alpha_z}{i\omega_c} = \frac{(\alpha_y + i\alpha_z)(-i)}{{\omega_c}} = \frac{-i\alpha_y + \alpha_z}{\omega_c} = \frac{\alpha_z}{\omega_c} - i\frac{\alpha_y}{\omega_c}.
\]

Note that $\alpha_y/\omega_c = E_y c/B \cdot (q/m)/(q/(mc)) = cE_y/B$ and similarly $\alpha_z/\omega_c = cE_z/B$.

Define the drift velocities:
\[
v_{dy} = \frac{cE_z}{B}, \qquad v_{dz} = -\frac{cE_y}{B}.
\]

(In SI units without the factor of $c$: $v_{dy} = E_z/B$, $v_{dz} = -E_y/B$.)

Then:
\[
w(t) = (v_{dy} - iv_{dz})(1 - e^{-i\omega_c t}).
\]

Separating real and imaginary parts with $e^{-i\omega_c t} = \cos\omega_c t - i\sin\omega_c t$:
\begin{align}
\dot{y} &= v_{dy}(1-\cos\omega_c t) - v_{dz}\sin\omega_c t, \\
\dot{z} &= -v_{dz}(1-\cos\omega_c t) - v_{dy}\sin\omega_c t.
\end{align}

\subsection*{Drift velocity}

Taking the time average over one cyclotron period ($\langle\cos\omega_c t\rangle = \langle\sin\omega_c t\rangle = 0$):
\[
\boxed{\langle\dot{y}\rangle = v_{dy} = \frac{E_z}{B}}, \qquad \langle\dot{z}\rangle = v_{dz} = -\frac{E_y}{B}.
\]

This is the well-known $\vect{E}\times\vect{B}$ drift: $\vect{v}_{\text{drift}} = c\,\frac{\vect{E}\times\vect{B}}{B^2}$, which is perpendicular to both $\vect{E}$ and $\vect{B}$.

The component $E_z$ produces a drift in the $y$-direction with magnitude $E_z/B$ (in SI units).

\subsection*{Positions by integration}

Integrating $\dot{y}$ and $\dot{z}$ with $y(0)=z(0)=0$:
\begin{align}
y(t) &= v_{dy}\!\left(t - \frac{\sin\omega_c t}{\omega_c}\right) + \frac{v_{dz}}{\omega_c}(\cos\omega_c t - 1), \\
z(t) &= -v_{dz}\!\left(t - \frac{\sin\omega_c t}{\omega_c}\right) - \frac{v_{dy}}{\omega_c}(\cos\omega_c t - 1).
\end{align}

\subsection*{Special case $E_z = 0$}

When $E_z = 0$, $v_{dy} = 0$ and $v_{dz} = -E_y/B$, so:
\begin{align}
y(t) &= \frac{v_{dz}}{\omega_c}(\cos\omega_c t - 1), \\
z(t) &= -v_{dz}\!\left(t - \frac{\sin\omega_c t}{\omega_c}\right) = v_{dz}\!\left(\frac{\sin\omega_c t}{\omega_c} - t\right).
\end{align}

The $y$-motion is purely oscillatory (bounded), while the $z$-motion has a secular drift (linear in $t$) plus oscillations. Combined with the uniform $x$-motion ($x = x_0 t$), the particle spirals about a line parallel to the $z$-axis with a drift in the $z$-direction. \hfill $\blacksquare$

\end{document}
